\begin{figure}[H]
\centering
\begin{minipage}[t]{0.48\textwidth}
  \vspace{0pt}\centering
  \includegraphics[width=\linewidth]{figures/ch08/panels/figure8_3_a.pdf}
\end{minipage}\hfill
\begin{minipage}[t]{0.48\textwidth}
  \vspace{0pt}\centering
  \includegraphics[width=\linewidth]{figures/ch08/panels/figure8_3_b.pdf}
\end{minipage}

\vspace{0.8em}

\begin{minipage}[t]{0.48\textwidth}
  \vspace{0pt}\centering
  \includegraphics[width=\linewidth]{figures/ch08/panels/figure8_3_c.pdf}
\end{minipage}\hfill
\begin{minipage}[t]{0.48\textwidth}
  \vspace{0pt}
  \caption{\textbf{从能量到归一化密度。}
  低能量区域获得较大的未归一化权重 $\exp(-E_\theta(x))$；除以 $Z_\theta$ 后即可变成概率密度而不改变其形状。给所有能量加同一个常数，会按相同比例缩放权重与 $Z_\theta$，因此归一化后的密度保持不变。}
  \label{fig:ebm-energy-to-density}
\end{minipage}
\end{figure}
